% =========================================================
% Stern--Brocot Limit Theory
% =========================================================

\subsubsection{Stern--Brocot Limit Theory}
\label{sec:stern-brocot-limit-theory}

\begin{remark}
The Stern--Brocot tree does more than enumerate the positive rationals.
Its infinite paths encode the positive irrationals as limits of rational
sequences.  The four results below build this theory from the ground up,
each depending on the previous one.
\end{remark}

% ----- lem:sb-denom-growth -----
\begin{lemma}[Denominator growth on infinite Stern--Brocot paths]
\label{lem:sb-denom-growth}
Let $a_0/b_0, a_1/b_1, a_2/b_2, \ldots$ be the sequence of endpoint
fractions encountered along any infinite path in the Stern--Brocot tree.
Then the denominators are strictly increasing:
\[
b_0 < b_1 < b_2 < \cdots,
\]
and in particular $b_n \to \infty$.
\hyperref[prf:sb-denom-growth]{Go to proof.}
\end{lemma}

\begin{remark*}[Standard quantified statement]
\[
\forall n \in \mathbb{N},\quad b_n < b_{n+1}.
\qquad
\forall M > 0\;\; \exists N \in \mathbb{N}\;\; \forall n \geq N,\quad b_n > M.
\]
\end{remark*}

\begin{remark*}[Interpretation]
At each step in the tree the new fraction produced by the mediant has
denominator equal to the sum of the two parent denominators, which
strictly exceeds each of them.  A path that turns left or right at every
level therefore visits fractions with strictly increasing denominators,
and since the denominators are positive integers they must diverge to
infinity.  This is the key lemma that makes the limit theory work: it
guarantees that the Stern--Brocot intervals shrink to zero length.
\end{remark*}

\begin{remark*}[Dependencies]
\begin{itemize}
  \item \hyperref[thm:farey-neighbor-theorem]{Farey Neighbor Theorem}
  \item \hyperref[thm:induction-principle-for-peano-system]{Induction Principle for a Peano System}
\end{itemize}
\end{remark*}

% ----- lem:sb-length-zero -----
\begin{lemma}[Stern--Brocot interval lengths shrink to zero]
\label{lem:sb-length-zero}
Let $(I_n)_{n \geq 0}$ be the nested rational intervals generated by
the Stern--Brocot refinement process targeting a real number $x > 0$,
where $I_n = [a_n/b_n,\; c_n/d_n]$.  Then the interval lengths
$\ell_n = c_n/d_n - a_n/b_n$ satisfy
\[
\lim_{n \to \infty} \ell_n = 0.
\]
\hyperref[prf:sb-length-zero]{Go to proof.}
\end{lemma}

\begin{remark*}[Standard quantified statement]
\[
\forall \varepsilon > 0\;\;
\exists N \in \mathbb{N}\;\;
\forall n \geq N,\quad
\ell_n < \varepsilon.
\]
\end{remark*}

\begin{remark*}[Interpretation]
By the Farey gap formula, $\ell_n = 1/(b_n d_n)$.
By the Denominator Growth Lemma, both $b_n$ and $d_n$ tend to infinity.
Therefore $1/(b_n d_n) \to 0$.  The Archimedean property of $\mathbb{Q}$
converts this into an explicit $N$: given $\varepsilon > 0$, find $N$
such that $b_N d_N > 1/\varepsilon$.
\end{remark*}

\begin{remark*}[Dependencies]
\begin{itemize}
  \item \hyperref[thm:farey-gap]{Gap between Farey Neighbors}
  \item \hyperref[lem:sb-denom-growth]{Denominator Growth on Infinite SB Paths}
  \item \hyperref[thm:archimedean-q]{Archimedean Property of $\mathbb{Q}$}
\end{itemize}
\end{remark*}

% ----- prop:sb-unique-point -----
\begin{proposition}[Nested Stern--Brocot intervals contain a unique limit point]
\label{prop:sb-unique-point}
Let $(I_n)_{n \geq 0}$ be the nested sequence of closed rational intervals
generated by the Stern--Brocot refinement process targeting $x > 0$,
with $I_n = [a_n/b_n,\; c_n/d_n]$.  Then there exists a unique real
number $\xi \in \bigcap_{n=0}^{\infty} I_n$.
\hyperref[prf:sb-unique-point]{Go to proof.}
\end{proposition}

\begin{remark*}[Standard quantified statement]
\[
\exists!\, \xi \in \mathbb{R}\;\;
\forall n \in \mathbb{N},\quad
\frac{a_n}{b_n} \leq \xi \leq \frac{c_n}{d_n}.
\]
\end{remark*}

\begin{remark*}[Interpretation]
The left endpoints $(a_n/b_n)$ form a monotone increasing sequence bounded
above; the right endpoints $(c_n/d_n)$ form a monotone decreasing sequence
bounded below.  Both converge in $\mathbb{R}$ by the Monotone Convergence
Theorem, and since the interval lengths tend to zero, the two limits must
coincide.  Uniqueness follows immediately: any two points in
$\bigcap I_n$ differ by at most $\ell_n$ for every $n$, hence by zero.

This is where completeness of $\mathbb{R}$ enters: the Monotone Convergence
Theorem requires the least-upper-bound property, which $\mathbb{Q}$ does
not have.  This is precisely why $\xi$ need not be rational.
\end{remark*}

\begin{remark*}[Dependencies]
\begin{itemize}
  \item \hyperref[lem:sb-length-zero]{Stern--Brocot Interval Lengths Shrink to Zero}
\end{itemize}
\end{remark*}

% ----- cor:sb-path-limit -----
\begin{corollary}[Every infinite Stern--Brocot path converges to its target]
\label{cor:sb-path-limit}
Let $x > 0$ be a real number and let $(I_n)$ be the Stern--Brocot
refinement sequence targeting $x$.  The sequences of left and right
endpoints both converge to $x$:
\[
\lim_{n \to \infty} \frac{a_n}{b_n}
= \lim_{n \to \infty} \frac{c_n}{d_n}
= x.
\]
\hyperref[prf:sb-path-limit]{Go to proof.}
\end{corollary}

\begin{remark*}[Standard quantified statement]
\[
\forall \varepsilon > 0\;\;
\exists N \in \mathbb{N}\;\;
\forall n \geq N,\quad
\left|x - \frac{a_n}{b_n}\right| < \varepsilon
\;\land\;
\left|x - \frac{c_n}{d_n}\right| < \varepsilon.
\]
\end{remark*}

\begin{remark*}[Interpretation]
By construction, $x \in I_n$ for every $n$: the refinement rule always
replaces the endpoint that would push $x$ outside the interval, so $x$
is retained in $\bigcap I_n$.  But the proposition says this intersection
contains exactly one point, namely $\xi$.  Therefore $\xi = x$, and the
endpoint sequences converge to $x$.
\end{remark*}

\begin{remark*}[Dependencies]
\begin{itemize}
  \item \hyperref[prop:sb-unique-point]{Nested SB Intervals Contain a Unique Point}
\end{itemize}
\end{remark*}

% ----- cor:sb-limit-irrational -----
\begin{corollary}[Infinite Stern--Brocot paths converge to irrationals]
\label{cor:sb-limit-irrational}
A positive real number $x$ corresponds to a \emph{finite} path in the
Stern--Brocot tree if and only if $x$ is rational.  Equivalently, $x$
corresponds to an \emph{infinite} path if and only if $x$ is irrational.
\hyperref[prf:sb-limit-irrational]{Go to proof.}
\end{corollary}

\begin{remark*}[Standard quantified statement]
\[
x \in \mathbb{Q}_{>0}
\;\Longleftrightarrow\;
\text{the SB refinement process for } x \text{ terminates in finitely many steps}.
\]
\end{remark*}

\begin{remark*}[Interpretation]
The refinement process terminates at the first step where $x$ equals the
mediant.  Since every mediant is rational, termination forces $x \in \mathbb{Q}$.
Conversely, if $x = p/q$ is rational, the Denominator Growth Lemma forces
the denominators past $q$ in finitely many steps, at which point $x$ must
have already been identified as a mediant.  So the rational numbers are
exactly those reached by finite paths, and the irrationals are exactly those
requiring infinite paths.  This is the cleanest formulation of the
incompleteness of $\mathbb{Q}$ via the Stern--Brocot tree.
\end{remark*}

\begin{remark*}[Dependencies]
\begin{itemize}
  \item \hyperref[thm:farey-neighbor-theorem]{Farey Neighbor Theorem}
  \item \hyperref[lem:sb-denom-growth]{Denominator Growth on Infinite SB Paths}
  \item \hyperref[thm:sqrt2-irrational]{No Rational Number Has Square $2$}
\end{itemize}
\end{remark*}
